\section{Definitions}

\subsection{Geometry}

\begin{definition}[Riemannian Metric and Volume Form]
  A \emph{Riemannian metric} $g$ on $M$ is a smooth symmetric positive-definite $(0,2)$-tensor field on $M$. In local coordinates $\{x^i\}$, $g$ is given by $g = g_{ij} \dd{x^i} \otimes \dd{x^j}$. The associated \emph{volume form} is
  \[%
    \dd{V}_g \coloneqq \sqrt{\det(g_{ij})} \, \dd{x^1} \wedge \cdots \wedge \dd{x^n}
  .\]%
\end{definition}

\begin{definition}[Covariant Derivative]
  The \emph{covariant derivative} $\nabla$ is the unique connection on $TM$ compatible with the metric $g$ and torsion-free, i.e.\ satisfying
  \[%
    \nabla g = 0, \qquad \nabla_X Y - \nabla_Y X = [X, Y]
  ,\]%
  for all vector fields $X, Y$ on $M$. For a smooth function $u \in C^\infty(M)$, the \emph{$i$-th order covariant derivative} $\nabla^i u$ is defined inductively by $\nabla^1 u = \nabla u$ and $\nabla^i u = \nabla(\nabla^{i-1} u)$.
\end{definition}

\begin{definition}[Geodesic Distance]
  The \emph{geodesic distance} $d(x, y)$ between two points $x, y \in M$ is the infimum of the lengths of all piecewise smooth curves $\gamma : [0, 1] \to M$ with $\gamma(0) = x$ and $\gamma(1) = y$, where the length of $\gamma$ is given by
  \[%
    \ell(\gamma) \coloneqq \int_0^1 \sqrt{g(\dot\gamma(t), \dot\gamma(t))} \dd{t}
  .\]%
\end{definition}

\begin{definition}[Isometry Group]
  The \emph{isometry group} of a Riemannian manifold $(M, g)$ is the group of all diffeomorphisms $\phi : M \to M$ such that $\phi^*g = g$, where $\phi^*g$ is the pullback of $g$ by $\phi$.
\end{definition}

\begin{definition}[Partition of Unity]
  A \emph{partition of unity} subordinate to an open cover $\{U_i\}$ of $M$ is a collection of smooth functions $\{\psi_i\}$ on $M$ such that
  \begin{enumerate}
    \item $0 \le \psi_i \le 1$ for all $i$,
    \item $\supp(\psi_i) \subseteq U_i$ for all $i$ (\cref{def:compactly_supported_functions}),
    \item $\sum_i \psi_i = 1$ everywhere on $M$.
  \end{enumerate}
\end{definition}

\subsection{Function Spaces}

\begin{definition}[Locally Integrable Functions]
  A function $u : M \to \R$ is \emph{locally integrable} if for every compact set $K \Subset M$,
  \[%
    \int_K |u| \dd{V}_g < \infty
  .\]%
  The space of locally integrable functions on $M$ is denoted $L^1_{\text{loc}}(M)$.
\end{definition}

\begin{definition}[Compactly Supported Functions]
  \label{def:compactly_supported_functions}

  A function $u : M \to \R$ has \emph{compact support} if the set $\supp(u) \coloneqq \overline{\{x \in M : u(x) \neq 0\}}$ is compact. We denote by $C^\infty_c(M)$ the space of smooth functions on $M$ with compact support.
\end{definition}

\begin{definition}[Lebesgue Space]
  If $q \ge 1$, the \emph{Lebesgue Space} $L^q(M)$ is the set of locally integrable functions $u$ on $M$ such that
  \begin{equation}\label{eq:lebesgue_norm}
    \|u\|_q \coloneqq \left(\int_M |u|^q \dd{V}_g\right)^{1/q}
  ,\end{equation}
  is finite.
\end{definition}

\begin{definition}[Sobolev Space]
  For $q \ge 1$ and a nonnegative integer $k$, the \emph{Sobolev Space} $L^q_k(M)$ is the set of $u \in L^q(M)$ whose covariant derivatives $\nabla^i u$ exist weakly and satisfy $\nabla^i u \in L^q(M)$ for all $0 \le i \le k$. The \emph{Sobolev norm} on $L^q_k(M)$ is given by
  \begin{equation}\label{eq:sobolev_norm}
    \|u\|_{q,k} \coloneqq \left(\sum_{i=0}^k \int_M \left\lvert\nabla^i u\right\rvert^q \dd{V}_g\right)^{1/q}
  .\end{equation}
\end{definition}

\begin{definition}[Space of $k$-times Differentiable Functions]
  \emph{The space of $k$-times differentiable functions}, $C^k(M)$, is the set of functions on $M$ that are $k$-times continuously differentiable. The \emph{norm} on $C^k(M)$ is given by:
  \begin{equation}\label{eq:ck_norm}
    \|u\|_{C^k} \coloneqq \sum_{i=0}^k \sup_M \left\lvert\nabla^i u\right\rvert
  ,\end{equation}
  is finite.
\end{definition}

\begin{note}
  Notice that $C^\infty(M) = \bigcap_{k=0}^\infty C^k(M)$, and $C^\infty(M) \subset L^q_k(M)$ for all $q$ and $k$.
\end{note}

\begin{note}
  We denote the space of continuous functions as $C^0(M) = C(M)$, and the norm on $C(M)$ is given by $\|u\|_{C^0} = \sup_M |u|$.
\end{note}

\begin{definition}[H\"older Space]
  The \emph{H\"older Space}, $C^{k,\alpha}(M)$, is defined for $0 < \alpha < 1$ as the set of $u \in C^k(M)$ for which the \emph{H\"older norm}
  \begin{equation}\label{eq:holder_norm}
    \|u\|_{C^{k,\alpha}} \coloneqq \|u\|_{C^k} + \sup_{\substack{x,y \in M \\ x \neq y}} \frac{\left\lvert \nabla^k u(x) - \nabla^k u(y)\right\rvert}{d(x, y)^\alpha}
  ,\end{equation}
  is finite, where $d(x, y)$ denotes the geodesic distance between $x$ and $y$.
\end{definition}

\subsection{Operators}

\begin{definition}[Laplacian]
  The \emph{Laplacian} $\Delta$ is the second-order differential operator defined by $\Delta u = \divg(\grad u)$, where $\divg$ is the divergence operator and $\grad$ is the gradient operator.
\end{definition}

\begin{definition}[Schr\"odinger Operator]
  The \emph{Schr\"odinger operator} is the second-order differential operator defined by $\Delta + h$, where $\Delta$ is the Laplacian and $h$ is a smooth function on $M$.
\end{definition}

\begin{definition}[Elliptic Operator]
  A second-order differential operator $L = \sum_{i,j} a^{ij}(x) \nabla_i \nabla_j + \text{lower order terms}$ is \emph{elliptic} at $x \in M$ if the symbol matrix $(a^{ij}(x))$ is positive definite, i.e.\ there exists $\lambda > 0$ such that
  \[%
    \sum_{i,j} a^{ij}(x) \xi_i \xi_j \ge \lambda |\xi|^2
  ,\]%
  for all $\xi \in \R^n$. The operator is \emph{elliptic on $M$} if it is elliptic at every point of $M$.
\end{definition}

\subsection{Analysis}

\begin{definition}[Weak Derivative]
  Let $u \in L^1_{\text{loc}}(M)$. A function $v \in L^1_{\text{loc}}(M)$ is the \emph{weak derivative} of $u$ of order $i$ if for all $\phi \in C^\infty_c(M)$,
  \[%
    \int_M v \, \phi \dd{V}_g = (-1)^i \int_M u \, \nabla^i \phi \dd{V}_g
  .\]%
  We write $\nabla^i u \coloneqq v$.
\end{definition}

\begin{definition}[Weak-Solution]
  A function $u \in L^1_{\text{loc}}(M)$ is a \emph{weak solution} to the equation $\Delta u = f$ if for all $\phi \in C^\infty_c(M)$,
  \[%
    \int_M u \, \Delta\phi \dd{V}_g = \int_M f \phi \dd{V}_g
  .\]%
\end{definition}

\begin{definition}[Green Function]
  A \emph{Green function} for the operator $\Delta + h$ is a function $G : M \times M \to \R$ such that for each fixed $y \in M$, $G(\cdot, y)$ is a weak solution to $(\Delta + h)G(\cdot, y) = \Delta_y$, where $\Delta_y$ is the Dirac delta distribution at $y$.
\end{definition}

\begin{definition}[Lipschitz Continuity]
  A function $u : U \to \R$ is called \emph{Lipschitz continuous} if
  \[%
    |u(x) - u(y)| \le C|x - y|
  ,\]%
  for some constant $C$ and all $x, y \in U$. We write:
  \[%
    \Lip[u] := \sup_{x,y \in U \atop x \neq y} \frac{|u(x) - u(y)|}{|x - y|}
  .\]%
\end{definition}

\subsection{Partial Differential Equations}

\begin{definition}[General PDE Form]
  An expression of the form:
  \begin{equation}\label{eq:general_pde_form}
    F(D^ku(x), D^{k-1}u(x), \cdots, Du(x), u(x), x) = 0
  ,\end{equation}
  is called a $k^{\text{th}}$-order partial differential equation (PDE), where:
  \[%
    F : \R^{n^k} \times \R^{n^{k-1}} \times \cdots \times R^n \times \R \times U \to \R
  ,\]%
  is given and $u : U \to \R$ is the unknown function we want to solve for.
\end{definition}

\begin{definitions}\leavevmode
  \begin{enumerate}
    \item The PDE from \cref{eq:general_pde_form} is called \emph{linear} if it has the form:
      \[%
        \sum_{|\alpha|\le k} a_\alpha(x) D^\alpha u(x) = f(x)
      ,\]%
      for given functions $a_\alpha$ and $f$. It's called \emph{homogeneous} if $f \equiv 0$.
    \item The PDE from~\cref{eq:general_pde_form} is called \emph{semilinear} if it has the form:
      \[%
        \sum_{|\alpha|=k} a_\alpha(x) D^\alpha u + a_0(D^{k-1}u, \cdots, Du, u, x) = 0
      .\]%
    \item The PDE from~\cref{eq:general_pde_form} is called \emph{quasilinear} if it has the form:
      \[%
        \sum_{|\alpha|=k} a_\alpha(D^{k-1}u, \cdots, Du, u, x)D^\alpha u + a_0(D^{k-1}u, \cdots, Du, u, x) = 0
      .\]%
    \item The PDE from~\cref{eq:general_pde_form} is called \emph{fully nonlinear} if it depends nonlinearly upon the highest order derivatives.
  \end{enumerate}
\end{definitions}

\begin{definition}[Convex]
  A function $f : \R^n \to \R$ is called \emph{convex} provided that
  \[%
    f(\tau x + (1 - \tau)y) \le \tau f(x) + (1 - \tau)f(y)
  ,\]%
  for all $x, y \in \R^n$ and each $0 \le \tau \le 1$.
\end{definition}

